% Generated draft: requires mathematical and source review.
\documentclass[11pt]{article}
\usepackage[margin=25mm]{geometry}
\usepackage{fontspec}
\setmainfont{Latin Modern Roman}
\usepackage{amsmath,amssymb,mathtools,unicode-math}
\setmathfont{Latin Modern Math}
\usepackage{enumitem,needspace,xurl,hyperref}
\setlength{\parindent}{0pt}
\setlength{\parskip}{5pt}
\setlength{\emergencystretch}{3em}
\newcommand{\N}{\mathbb{N}}
\newcommand{\Z}{\mathbb{Z}}
\newcommand{\Q}{\mathbb{Q}}
\newcommand{\R}{\mathbb{R}}
\newcommand{\C}{\mathbb{C}}
\newcommand{\F}{\mathbb{F}}
\newcommand{\A}{\mathbb{A}}
\newcommand{\PP}{\mathbb{P}}
\newcommand{\E}{\mathbb{E}}
\newcommand{\eps}{\varepsilon}
\newcommand{\dd}{\,\mathrm{d}}
\begin{document}
\section*{UnsolvedMath ID: 20000413}
\textbf{Source problem code:} catalogue entry
\subsection{Definitions and mathematical statement}
K. Murty's AIM question starts with an automorphic form \(\pi\) over \(\Q\), a real quadratic field \(F\), and its base change \(\Pi\). It describes cohomological summands \(H^1(\pi)\subset H^1(M)\) for a modular curve and \(H^2(\Pi)\subset H^2(S_F)\) for a Hilbert modular surface, with
\[
H^2(\Pi)|_{G_F}\cong H^1(\pi)\otimes H^1(\pi).
\]
For two such fields \(F_1,F_2\), put \(L=F_1F_2\). The resulting summand of \(H^4(S_{F_1}\times S_{F_2})\) restricts to \(V^{\otimes4}\), \(V=H^1(\pi)\), as a \(G_L=\operatorname{Gal}(\overline\Q/L)\)-module. The question asks for the expected codimension-two algebraic cycles defined over \(L\). Precise cohomology, coefficient, compactification, automorphic weight, and character conventions are not specified; actual Tate classes belong to \(H^4(\,\cdot\,,\Q_\ell(2))^{G_L}\), not the untwisted full tensor space.

The attached report imposes an explicit big-monodromy hypothesis: over a characteristic-zero coefficient field, \(V\) is two-dimensional and the Zariski closure of the \(G_L\)-image contains \(SL(V)\). Put \(D=\det V\). It asserts
\[
\dim\left(V^{\otimes4}\otimes D^{-2}\right)^{G_L}=2.
\]
With alternating-pair tensors \(t_{ij\mid kl}\) in the indicated tensor positions and determinant normalization, its stated basis is
\[
p=t_{12\mid34},\qquad c=t_{13\mid24}+t_{14\mid23}.
\]
The first line is the external product of the two determinant/divisor lines. The second is the unique non-product Casimir line in
\(\operatorname{Ad}^0(V)\otimes\operatorname{Ad}^0(V)\), where \(\operatorname{Ad}^0(V)\) is the trace-free endomorphism representation, and corresponds via its invariant pairing to the identity on \(\operatorname{Ad}^0(V)\). Tensor sign and normalization conventions are not fully supplied.

After restoring the actual codimension-two Tate twist, the report gives the global characters of these two lines as
\[
\epsilon^2\chi_1\chi_2\quad\text{and}\quad\epsilon^2,
\]
respectively. Here \(\chi_i\) is the quadratic character associated to \(F_i/\Q\); the precise convention for \(\epsilon\) must be fixed through the determinant of the automorphic representation. In compatible notation it is the finite character left by \(D(1)\). Since \(\chi_i|_{G_L}=1\), the actual classes over \(L\) require the corresponding restriction of \(\epsilon^2\) to be trivial. Two determinant-normalized invariant lines do not by themselves ensure two untwisted or incorrectly twisted Tate classes. The source does not state the needed central-character restriction.

The report reduces the non-product geometric target to constructing an algebraic adjoint identity correspondence between the two surface summands. It explicitly does not construct that codimension-two cycle for distinct \(F_1,F_2\). Thus the two-line computation is a representation-theoretic reduction, not a proof of algebraicity of the desired non-product class.
\subsection{Short English statement}
Find codimension-two cycles linking Hilbert modular surfaces from two quadratic base changes. Under big monodromy the report isolates two determinant-normalized lines and their descent characters, leaving the non-product adjoint correspondence unconstructed.
\subsection{Sources}
\begin{itemize}
\item AIM workshop ``The Tate conjecture,'' Question 4, attributed to K. Murty: \url{https://aimath.org/WWN/tateconjecture/tateconjecture.pdf}. No PDF was accessed here.
\item Local canonical registry \texttt{aim-algebraic-number-theory-notes.json}, \texttt{notes[64]}, tag \texttt{problem}; permanent UnsolvedMath ID \(20000413\), source code \texttt{catalogue entry}. Original describes the two base changes and tensor-fourth-power summand. The verification identifies the final isolated 1 as a page number and restores lost bars and superscripts. Original remarks and literature field are empty.
\item Supplied research attempt 1 proposing the big-monodromy invariant calculation, explicit tensor basis, and descent-character audit. No authoritative full bibliographic references or geometric cycle construction are supplied. The report is machine generated, not peer reviewed, and requires independent verification before citation.
\end{itemize}
\end{document}
